\chapter{Methodology}
\label{chap:methodology}

% ============ GUIDELINE AUDIT 2026-07-25 -- METHODOLOGY (50% CRITERION) ============
% "Technical content and project execution" carries 50% -- half the dissertation's mark. The 70+
% band reads: "Appropriate choice of strategy and approach. Clear and logical design. Detailed
% explanation showing evidence of reading and understanding of methods. METHODOLOGY EXECUTED
% RIGOROUSLY. Limitations recognised. An ambitious project executed well. Work at the higher end
% of this category may be suitable for publication in a conference or journal venue."
% Full audit: docs/TCD_Dissertation/GUIDELINE_COMPLIANCE.md
%
% The design work here is genuinely strong and already clears most of the 70+ band: the staged
% single-factor ladder, the dual-lens judge, the deterministic probes kept as a judge-free anchor,
% the tool-profile control, and the decoupling of generation from judging are all rigour the band
% asks for. Four gaps stand between this and the TOP of that band, all of which are fixable from
% artefacts already saved -- no GPU time, no new runs.
%
% (A) THE JUDGE MODEL IS NEVER NAMED. Section~\ref{subsec:mr-judge-validity} says only "a more
%     capable language model" and Section~\ref{subsec:mr-runbook} says "frontier models served
%     through NVIDIA NIM". A study whose headline numbers come from an unnamed grader is not
%     reproducible, and reproducibility is the load-bearing claim of the Runbook subsection. The
%     [AUDIT 2026-07-24] block at the end of this chapter already records the actual judges and
%     notes the NIM claim is inaccurate for the comparative lens. Apply it.
%
% (B) NO UNCERTAINTY IS REPORTED ON ANY NUMBER. This is the single biggest evidence gap in the
%     dissertation. The headline finding of the third hypothesis is a PARITY claim -- 0.589 against
%     0.583, described in Chapter~\ref{chap:experiments-results} as "within the noise of a
%     single-judge protocol". No interval, no test, no per-cell n. Section~\ref{subsec:evaluation-limitations}
%     concedes "a modest number of test items yields wide uncertainty" in words but never gives the
%     number. Under "evidence of critical and WELL SUPPORTED analysis", a stated-but-unquantified
%     uncertainty reads as an unsupported assertion -- and "presents some unsupported assertions"
%     is verbatim the 50-60 band descriptor on the presentation criterion.
%     FIX (cheap, no GPU): the 151 questions are paired across all five conditions and every verdict
%     is already persisted. A paired bootstrap over questions yields a confidence interval for each
%     H2H score and each ladder delta directly from the saved reports. ACTION: add an interval
%     column to tab_rank_overall and tab_rank_lineage, state the per-suite n beside each score, and
%     replace the phrase "within the noise of a single-judge protocol" with the measured interval.
%     Add one sentence to Section~\ref{subsec:mr-suites-metrics} stating the resampling procedure.
%
% (C) AN ABLATION THAT WAS RUN IS NEVER REPORTED. Appendix~\ref{chap:judge-prompts} states that the
%     absolute judge was re-run under two further exposures -- "bare" (principle name only, no
%     rubric) and "none" (no principle at all) -- as the self-enhancement control. The RESULTS of
%     that ablation appear nowhere in the body. This is completed validation work, and it is
%     precisely the evidence that answers the obvious examiner question, "is the judge rewarding
%     constitution-shaped text rather than actual compliance?". Leaving it unreported gives away
%     rigour that has already been earned. ACTION: add a short table to
%     Section~\ref{subsec:mr-judge-validity} giving each condition's score under full / bare / none
%     exposure, and one sentence on what the gap between them shows.
%
% (D) JUDGE REPEAT VARIANCE IS COLLECTED BUT NOT REPORTED. Section~\ref{subsec:mr-judge-validity}
%     states each item "is judged several times and averaged". The spread across those repeats is a
%     free reliability statistic and it is never given. ACTION: state the number of repeats and the
%     mean within-item standard deviation in one sentence. Together with (B) this converts the
%     judge from "a design argument" into a characterised instrument -- which is exactly the wording
%     Section~\ref{subsec:evaluation-limitations} uses to concede the current shortfall.
%
% NOTE on what does NOT need fixing: the absence of a human-agreement study is correctly scoped as
% future work and honestly conceded, and the guideline explicitly values that ("Do not be afraid to
% discuss disadvantages of your approach. Negative results are useful."). Keep that framing.
% ==================================================================================

A claim about trustworthy behaviour is only as sound as the instrument that measures it, so this chapter sets out that instrument before any result is reported. It describes the single apparatus inside which all three experiments run --- one constitution, one training corpus, one set of deterministic probes and language-model judges, and five model conditions, each held fixed so that a later experiment can change exactly one ingredient --- and, with it, what the on-device model of Figure~\ref{fig:architecture} is taught, what it is served, and how its behaviour is checked.

This chapter sets out the methodology of the dissertation: how the framework was built and how its behaviour was measured, so that the experiments in Chapter~\ref{chap:experiments-results} can be read as tests run inside one fixed environment, the ``apparatus'', rather than as a set of separate studies. Section~\ref{sec:mr-overview} explains the design that every experiment runs inside and the reasoning behind it. Section~\ref{sec:mr-shared-apparatus} then introduces the apparatus shared by all three experiments: the constitution being taught, the data used to teach it, the probes and judges used to measure behaviour, the model conditions being compared, and the infrastructure the runs were executed on. Each experiment's own method and results are then presented together in Chapter~\ref{chap:experiments-results}.

% [AUDIT 2026-07-24 -- framing / narrative through-line. Owen (supervisor) asked that the write-up
% ACTION: CONSIDER (judgement call, not mechanical): preface the chapter with one short paragraph (or a figure) stating the WHY/WHAT/HOW/RESULTS arc set out below.
% carry a clear message: a trustworthy, privacy-preserving, on-device framework -- the reader should
% know WHY and WHAT before HOW and RESULTS. The chapter currently opens straight into "the apparatus".
% Consider prefacing with one short paragraph (or a figure) that states the arc explicitly:
%   WHY  -- personal data must not leave the device (GDPR Art. 9 special-category data; Section~
%           \ref{subsec:gdpr-and-the-case-for-local-inference}), which forces a sub-1B on-device model.
%   WHAT -- a single reproducible framework/"apparatus" that turns a written constitution into a trained
%           on-device assistant AND a defensible evaluation of it (not five separate studies).
%   HOW  -- one pipeline of six stages: (1) generate culturally-diverse questions, (2) distil ideal
%           constitutional trajectories from a frontier teacher, (3) LoRA-train the 0.6B student,
%           (4) serve with a live tool loop, (5) benchmark across five suites, (6) judge with anchored
%           rubrics. (This is the six-stage table/mermaid in wiki/sources/code/pipeline-overview.md --
%           a full end-to-end pipeline FIGURE here (currently only the 4-stage DATA figure exists,
%           Figure~\ref{fig:data-pipeline}) would land "this is how it is done" at a glance.)
%   RESULTS -> forward-point to Chapter~\ref{chap:experiments-results}.
% This preface would also directly answer the "why / what / how / results" message the supervisor wants.
\section{Overview and Design Rationale}
\label{sec:mr-overview}

The three experiments follow a staged experimental design (Figure~\ref{fig:ladder}): each introduces one additional component over its predecessor, so that the effect of that component can be attributed to it alone.

\begin{figure}[H]
  \centering
  \begin{tikzpicture}[>=Stealth,
      stage/.style={draw, rounded corners, fill=black!4, align=center,
        text width=3.2cm, minimum height=2.2cm, inner sep=5pt, font=\footnotesize},
      arrowlab/.style={align=center, font=\scriptsize\itshape, text width=2.3cm}]
    \node[stage] (e1) {\textbf{Experiment 1}\\[2pt] template SFT\\ single 0.6B model\\ (format, not content)};
    \node[stage, right=2.5cm of e1] (e2) {\textbf{Experiment 2}\\[2pt] constitutional distillation\\ single 0.6B model\\ (frontier-teacher data)};
    \node[stage, right=2.5cm of e2] (e3) {\textbf{Experiment 3}\\[2pt] Thinker--Executor\\ two 0.6B models\\ (same training data)};
    \draw[->, thick] (e1) -- node[above, arrowlab] {format alone helps nothing: data or scale?} (e2);
    \draw[->, thick] (e2) -- node[above, arrowlab] {compliance up, reasoning not: divide the work} (e3);
  \end{tikzpicture}
  \caption{The staged experimental design: each experiment is the structural response to the previous one's measured result, read against the two untrained baselines of Section~\ref{subsec:mr-conditions}.}
  \label{fig:ladder}
\end{figure}

The design is organised and based around five requirements the dissertation chose to explore. The first is \emph{privacy by on-device execution}, when an assistant answers, no part of the user's conversation should be sent to the cloud, which rules out any design that leans on a remote cloud model. The second is the \emph{small model}, where the model must be small enough to run on a phone. The third is \emph{single-factor attribution}, where each experimental stage changes exactly one ingredient over its predecessor, so that any measured difference can be credited to that ingredient alone. The fourth is \emph{reproducibility}, the experiments must be repeatable from saved artefacts rather than re-created by hand. The fifth is \emph{deterministic, auditable checking}, the primary measure of whether the model behaved well must return the same verdict on the same input every time, so that a result can be trusted and re-run.

\section{Shared Apparatus}
\label{sec:mr-shared-apparatus}

To meet these requirements, all three experiments share one pipeline, the ``apparatus'' of this dissertation: a single pipeline for training, inference and benchmarking, so that external factors stay fixed while only the intervention changes. Its components are introduced in turn in the subsections below.

\subsection{The Constitution: Principles and Families}
\label{subsec:mr-constitution}

To frame the definition of ``trustworthy'' from Section~\ref{sec:definitions} the dissertation introduces a written constitution of twenty-five principles (P1--P25). The principles take inspiration from Anthropic's Constitutional AI framework \cite{bai2022constitutionalaiharmlessnessai} and the published Claude constitution \cite{anthropic2023claudesconstitution}, and are extended with domain-specific rules for personalisation and well-being; their full statements are given in Appendix~\ref{chap:constitutional-principles-reference} and the document itself is in the repository (\texttt{pipeline/constitution.md}).

% [DRAFT 2026-07-25 | background/LR (15%) + testing (15%) | HIGH VALUE] MAKE THIS A PREDICTION.
% This is the strongest dot-connecting opportunity in the dissertation and it is currently being
% wasted. The C3AI typology says prohibitions instil more readily than generative instructions.
% Your results then show exactly that -- robustness and honesty improve, the generative reasoning
% principles lag. But as the document stands, C3AI is introduced HERE as a neutral grouping rationale
% and only cashed in at Section~\ref{subsec:h1-constitutional-distillation}, where it reads as a
% post-hoc explanation found after the fact.
%
% A prediction stated BEFORE the result and confirmed after is the single clearest demonstration of
% "critical and analytical perspective" and "ability to relate theoretical knowledge to project work"
% -- the exact wording of the 70+ literature-review band. It costs one sentence here and three words
% in the results.
%
% (1) APPEND to the paragraph below:
%     "This typology also yields a prediction that the results can test rather than merely a
%      convenient grouping: if prohibitions are indeed easier to instil than generative instructions,
%      then constitutional fine-tuning should improve the negatively framed robustness and honesty
%      families more than the positively framed reasoning and personalisation families. The families
%      are therefore scored separately throughout, so that the prediction can be checked against the
%      per-family results of Chapter~\ref{chap:experiments-results} rather than assumed."
%
% (2) THEN, in Section~\ref{subsec:mr-exp2-results} where the per-family results are first reported,
%     REPLACE the neutral phrasing with an explicit callback, e.g. prepend to that sentence:
%     "As the C3AI typology predicted (Section~\ref{subsec:mr-constitution}), ..."
%
% (3) The Discussion paragraph at Section~\ref{subsec:h1-constitutional-distillation} can then be
%     SHORTENED, because it no longer has to introduce the idea -- it only has to confirm it. Its
%     current closing ("so the unevenness measured here is the predicted unevenness, and not a defect
%     in the method") becomes much stronger once the prediction was genuinely made in advance.
%
% This same before/after treatment is available for two other claims already in the document and is
% worth applying to both: LIMA predicts a small curated corpus will suffice (confirmed by
% Experiment~2's corpus size), and Zhang's collapse finding predicts self-critique will not transfer
% downward (which is why the student never critiques itself). State each as a prediction where the
% literature is introduced, then confirm it by name in the results.
Rather than treating each and every principle from the constitution as a flat checklist, they are grouped into five families, which are summarised in Table~\ref{tab:constitution}, because each family teaches a different kind of behaviour. The grouping follows the C3AI typology of Kyrychenko et al.\ \cite{kyrychenko2025c3aicraftingevaluatingconstitutions}, which sorts principles by how they are phrased, prohibitions to be avoided versus generative instructions to be carried out.

\begin{table}[t]
  \centering
  \caption{The five constitutional principle families and their scored items (twenty-one scored with P2/P3 merged; MEM is the memory probe; P22--P25 are defined but unprobed).}
  \label{tab:constitution}
  \begin{tabular}{lp{4cm}p{5cm}}
    \hline
    Family & Scored items & Framing \\
    \hline
    Reasoning and process & P1, P15, P20, P21 & positively framed behaviours \\
    Honesty and calibration & P5, P7, P8, P16, P18 & prohibitions and acknowledgements \\
    Tool discipline and use & P2/P3, P4, P10, P11, P12, P13, P19 & instrumental behaviours \\
    Robustness and integrity & P9, P14 & negatively framed traits \\
    Context and personalisation & P6, P17, MEM & positively framed behaviours \\
    \hline
  \end{tabular}
\end{table}

Principles are additionally weighted into three a-priori priority tiers. The tiers follow from the model's stated purpose (Chapter~\ref{chap:introduction}): ask the right question, never give a false answer, deny when it does not know, hold trust under pressure, personalise, and use tools to scrutinise its own thinking. Tier~1 ($\times 3$) holds the trust-critical \emph{outcomes} (ask the right question, no fabrication, deny-when-unknown), where a single failure most damages trust. Next, Tier~2 ($\times 2$) holds the \emph{substrate} that produces and sustains those outcomes, rigorous self-correcting reasoning (decompose, first-principles, self-correction), robustness under pressure, and the personalisation and memory substrate. Finally, Tier~3 ($\times 1$) holds the instrumental tool mechanism through which the work is carried out. Table~\ref{tab:tiers} maps each purpose onto its principles and weight, and the weighted score is $\sum_k w_k s_k / \sum_k w_k$ over the scored items $k$. The weighting therefore rewards obeying the principles that matter most, not obeying the greatest number of them.

However, asking the right question does not mean asking more: the model should pose a targeted 5W+H question only when the needed context is genuinely missing, and must not clarify when the user model already answers it; over-clarification is itself a failure. Avoiding false answers is judged with zero tolerance: a single fabrication fails the item outright and dominates the persona trustworthiness dimension rather than being averaged away. Denying when unknown requires an explicit acknowledgement of the gap (a knowledge cutoff, the absence of live data, or a plain ``I don't know''), not a vague hedge. Tool use is credited only when it is appropriate to the task. Personalisation counts only when retained facts are used accurately and remain correctable and a need for that question; fabricating a user fact is a trust failure, not a personalisation win.

\begin{table}[H]
  \centering
  \caption{A-priori purpose weighting of the constitution, fixed from the model's purpose rather than the results and always reported beside the unweighted mean.}
  \label{tab:tiers}
  \begin{tabular}{p{5cm}p{6.4cm}c}
    \hline
    Purpose & Scored items & Weight \\
    \hline
    Trust outcomes: ask the right question; never fabricate; deny when unknown
      & P21, P17, P6;\ P13, P5, P8;\ P18, P16, P7 & $\times 3$ \\
    Substrate: reason rigorously; hold trust under pressure; personalise
      & P1, P20, P15;\ P9, P14;\ MEM (memory) & $\times 2$ \\
    Instrumental tool mechanism
      & P4, P10, P12, P19, P2/P3, P11 & $\times 1$ \\
    \hline
  \end{tabular}
\end{table}

\subsection{Dataset Construction}
\label{subsec:mr-dataset}

Every fine-tuned model in this dissertation is produced by distillation from worked examples rather than from hand-written targets, but in two deliberately different generations. The first, used for Experiment~1, is a \emph{template} generation: each example is assembled from a fixed scaffold that fixes the response format (a \texttt{<think>} block followed by an answer) without exposing the model to the reasoning behind a constitutional decision. Its thought process follows a structured checklist, which bounds the model to the examples it was trained on. The second, used for Experiments~2 and~3, teaches the behaviour the constitution describes rather than its shape. A large frontier teacher answers each question with the full constitution in front of it, and the student is then shown only the answer, never the rules the teacher was reading, so the student learns to act as the constitution asks without ever seeing it; this one-sided view, the teacher sees the constitution and the student does not, is what is meant by calling the distillation \emph{asymmetric}, an idea that carries Anthropic's Constitutional AI (Section~\ref{sec:constitutional-ai}) down to a model too small to reason over the rules for itself.

% [AUDIT 2026-07-24 -- coverage gap Owen flagged as a strength. The supervisor specifically
% ACTION: CONSIDER (judgement call, not mechanical): add two to three sentences (or a small table of the sixteen categories against sample diversity axes) making the question-generation breadth explicit, and state the exact question count for reproducibility, per the note below.
% highlighted that "the pipeline generates questions based on different contexts and nationalities
% using an LLM, resulting in around 2,500 questions". That breadth is a genuine methodological
% contribution and is currently compressed to "a question set is generated to span the twenty-five
% principles" -- the two axes that give the corpus its diversity are not stated at all. Per
% wiki/sources/code/pipeline-overview.md (Stage 1), question generation cycles through TWO axes:
%   (a) SITUATIONS -- 16 question categories, each tied to the principle it stress-tests
%       (user_context_behavioral, real_time_dependent, impossible_tasks, subjective_tradeoffs,
%        adversarial_pressure, knowledge_boundary, multi_step_clarification, ambiguous_underspecified,
%        entity_facts_web_search, verbose_context_behavioral, multi_turn_conversation,
%        interleaved_tool_reasoning, scratchpad_decomposition, partial_capability_honest,
%        inventory_constraint, environment_timeout). The last two are deliberately adversarial
%        ENVIRONMENTS -- inventory_constraint removes a needed tool (correct behaviour = notice & refuse),
%        environment_timeout returns HTTP 503 on first search (correct behaviour = retry-once-then-fallback)
%        -- this is how "know your limits" negatives enter the corpus, not only happy-path examples.
%   (b) PERSONALITIES / NATIONALITIES -- 20 region/culture/demographic diversity axes (South Asia, East
%       Africa, Southeast Asia, Latin America, Middle East, diaspora communities, ...), >=60% of each batch
%       forced to reflect the slot, with local currencies (Rs, NGN, PHP, XOF, KES), local tax regimes
%       (GST India, VAT EU, HST Canada) and local financial/social practice (halal finance, chit funds,
%       stokvel, M-Pesa, UPI). This is the "who is the user" variation the personalisation/empathy claims
%       later depend on, and the concrete answer to a reviewer's "is this Western-defaults data?".
% Formats also vary: single-turn, two_turn (ask + pushback for adversarial_pressure), multi_turn (3-5
% messages revealing context progressively), verbose_single_turn (a paragraph of personal context).
% RECOMMEND: add 2-3 sentences (or a small table of the 16 categories x sample axes) making this explicit,
% and state the exact QUESTION count for reproducibility (Owen said ~2,500; the trained CORPUS is 2,897
% examples -- clarify the question-vs-example distinction so the two numbers don't read as a conflict).
The generation proceeds in the four stages of Figure~\ref{fig:data-pipeline}. A question set spanning P1--P25, with GSM8K items added for multi-step arithmetic \cite{cobbe2021trainingverifierssolvemath}, is answered by a frontier teacher (\texttt{minimax-m2.7}, served through NVIDIA NIM \cite{nvidianim}) whose tool calls are executed live, so each trajectory records real Python execution and web-search results (via Exa \cite{exaai}) rather than the teacher's guess at them. The teacher's constitution-laden prompt is then replaced by a short student prompt listing only the tools, so the student sees compliant trajectories but never the twenty-five principles; a length-based quality gate and robustness variants complete the assembly. The result is 2,897 constitutional examples, size-matched to the 2,895 template examples of Experiment~1 so the two corpora differ in content, not volume; the full composition, including the Thinker and Executor views, is tabulated in Appendix~\ref{chap:training-data-statistics}.

\begin{figure}[H]
  \centering
  \begin{tikzpicture}[>=Stealth,
      stage/.style={draw, rounded corners, fill=black!4, align=center,
        text width=2.95cm, minimum height=2.3cm, inner sep=4pt, font=\scriptsize}]
    \node[stage] (q) {\textbf{1. Question set}\\[2pt] spans P1--P25;\\ GSM8K arithmetic items};
    \node[stage, right=0.55cm of q] (t) {\textbf{2. Frontier teacher}\\[2pt] \texttt{minimax-m2.7}; full constitution in its prompt; tool calls executed live (Python, exa.ai)};
    \node[stage, right=0.55cm of t] (s) {\textbf{3. Asymmetric swap}\\[2pt] teacher prompt removed; short student prompt (tools only) saved instead};
    \node[stage, right=0.55cm of s] (a) {\textbf{4. Assembly}\\[2pt] robustness variants; think-length quality gate};
    \draw[->, thick] (q) -- (t);
    \draw[->, thick] (t) -- (s);
    \draw[->, thick] (s) -- (a);
    \node[stage, below=0.9cm of a] (corp) {\textbf{Constitutional corpus}\\[2pt] 2,897 examples;\\ trains the single model (Experiment~2)};
    \node[stage, left=0.55cm of corp] (split) {\textbf{Trajectory splitter}\\[2pt] projects each trajectory into two role views (Experiment~3)};
    \node[stage, left=0.55cm of split] (views) {\textbf{Role views}\\[2pt] Thinker: 2,720 examples\\ Executor: 2,243 examples};
    \draw[->, thick] (a) -- (corp);
    \draw[->, thick] (corp) -- (split);
    \draw[->, thick] (split) -- (views);
  \end{tikzpicture}
  \caption{Construction of the training data: the constitution is shown to the teacher, not the student, and the same corpus feeds both the single model (Experiment~2) and the Thinker--Executor split (Experiment~3).}
  \label{fig:data-pipeline}
\end{figure}

% [DRAFT -- dataset samples for all three experiments; addresses "add a sample of the dataset and the
% output it published" across Exp 1--3. One characteristic training row per corpus makes the staged design
% concrete and lets the prose describing each generation be trimmed (a reader sees the contract rather than
% only reads it described). All rows are real and unedited, abridged with \ldots only for length:
% Template = train_interleaved_native.jsonl (2,895); Constitutional = train_sft_v3.jsonl (2,897; qid
% f80b77cf7a63); Thinker/Executor = train_sft_thinker_curriculum.jsonl (2,720) and train_sft_executor.jsonl
% (2,243; qid 14c514421935). To place it: uncomment and add a pointer such as "one characteristic row from
% each corpus is shown in Table~\ref{tab:dataset-sample}" to the paragraph above. Benchmark OUTPUTS are
% already shown in the case-study tables (tab:case-p5/tab:case-p6, Experiments chapter).
% \begin{table}[H]
%   \centering\small
%   \caption{One characteristic training row from each experiment's corpus (real, abridged), showing what
%     the student is trained on at each stage. Experiment~1 teaches a fixed response \emph{format};
%     Experiment~2 teaches constitutional \emph{behaviour} distilled from a teacher that read the
%     constitution the student never sees; Experiment~3 projects one constitutional trajectory onto two role
%     views whose actions match by construction.}
%   \label{tab:dataset-sample}
%   \begin{tabular}{@{}p{2.6cm}p{10.1cm}@{}}
%     \toprule
%     corpus & characteristic training row (abridged) \\
%     \midrule
%     \textbf{Template} \newline (Exp~1; 2,895) \newline \emph{format, not content}
%       & \texttt{user}: ``Recommend me a book.'' \quad \texttt{<think>}: ``UNDERSTAND: the user wants a book recommendation. PLAN / CLASSIFY / ANALYZE: genre UNKNOWN; recent reads UNKNOWN; goal UNKNOWN\ldots{} CHECKPOINT: have it? NO --- ASK USER.'' \quad \texttt{<answer>}: ``I'd love to recommend something\ldots{} 1.~Genre? 2.~Recent reads? 3.~Goal? 4.~Length?'' \newline \emph{A fixed checklist scaffold; the constitution is never in view.} \\
%     \addlinespace
%     \textbf{Constitutional} \newline (Exp~2; 2,897) \newline \emph{asymmetric distillation}
%       & \texttt{user}: ``Should I pursue a master's degree in data science in the UAE or study abroad in Europe?'' \quad \texttt{<think>}: ``This question depends entirely on who the user is\ldots{} reading user memory and personalising is the most important thing here.'' \quad \emph{memory read} $\rightarrow$ ``[WHO] bakery owner in Madrid\ldots{} [CONSTRAINTS] very limited time; tight budget.'' \quad \texttt{<answer>}: ``\ldots{}your stored profile\ldots{} sits quite far from the data science master's space\ldots{} is this decision for \textbf{you personally}, or\ldots{} someone else?'' \newline \emph{The teacher read the constitution; the student sees only the compliant trajectory.} \\
%     \addlinespace
%     \textbf{Thinker / Executor} \newline (Exp~3; 2,720 / 2,243) \newline \emph{one trajectory, two views}
%       & \emph{Thinker view.} \texttt{user}: ``\ldots{}a job offer in Gothenburg\ldots{} is the healthcare and education quality significantly better to justify the higher taxes?'' \quad \texttt{<think>}: ``\ldots{}the question is concrete enough to answer with current data\ldots{} a web search is the right next step.'' \quad \texttt{<act>}: ``Search for current comparisons of Swedish healthcare and education quality\ldots{}'' \newline \emph{Executor view (same trajectory).} \texttt{user}: ``Use web\_search to find: Sweden healthcare education quality compared high taxes'' \quad \texttt{<tool\_call>}: \texttt{web\_search(\{query\})}. \newline \emph{The \texttt{<act>} the Thinker requests and the call the Executor makes are the same action by construction.} \\
%     \bottomrule
%   \end{tabular}
% \end{table}

Experiment~3's data is created by splitting each conversation along the boundary between thinking and execution: a \emph{Thinker view} that keeps the reasoning and the \texttt{<think>}/\texttt{<ask>}/\texttt{<act>}/\texttt{<answer>} control flow, and an \emph{Executor view} that maps a single \texttt{<act>} instruction to a single tool call. Because both views are read from the same trajectory, the action the Thinker learns to request and the call the Executor learns to make are by construction the same action, so their alignment is structural rather than hoped for. The Executor owns the external-world tools (Python execution, web search, page reading), while the Thinker's own state tools (user memory, scratchpad) are supplied by the orchestrator at inference and are dropped from its training stream with their reasoning preserved. One consequence matters for Experiment~3's interpretation: the single constitutional model and the Thinker--Executor are trained on the same trajectories, so their comparison isolates the architecture and not the data.

\subsection{System Prompts at Training and Serving Time}
\label{subsec:mr-system-prompts}

Three families of system prompt run through this study, and keeping them apart is essential to understanding what each experiment tests. The first is the \emph{teacher prompt}, used only during data generation and never saved into a training example; it embeds the full written constitution as behavioural principles, together with strict format rules for the reasoning and tool-call structure of each trajectory. The second is the \emph{student prompt}, a compact instruction of roughly two hundred and thirty words that states the assistant's role, the required output structure (reason first in \texttt{<think>}, then answer in \texttt{<answer>}), a tool policy for the session, and a set of security rules; it never mentions the constitution. % [REPETITION -- helps trimming] "the principles enter the weights rather than the context / the student
% ACTION: CONSIDER (judgement call, not mechanical): keep this statement once where asymmetric distillation is first defined (Section~\ref{subsec:mr-dataset}) and shorten the restatement here to a back-reference.
% sees only compliant behaviour, never the rules" is stated three times: twice in Section~\ref{subsec:mr-dataset}
% ("distilled into the weights as behaviour rather than copied into the context as text"; "the student
% therefore never reads the twenty-five principles") and here. Keep it once where asymmetric distillation is
% first defined and shorten to a back-reference here.
The asymmetry between these two prompts is the mechanism of the distillation described in Section~\ref{subsec:mr-dataset}: the teacher is shown the rules while the student is shown only behaviour that already obeys them, so the principles enter the weights rather than the context.

The student prompt is also where training and serving meet. The very same prompt string that is saved into every training example is loaded by the inference server from the same source at serving time, so what the model was trained under and what it is served under cannot drift apart. Three variants exist, differing only in the single line that states which tools are available in the session (full tools, computation only, or no external tools), matching the benchmark's tool profiles.

The Thinker--Executor architecture adds two role prompts of its own, kept under the same single-source-of-truth discipline and shared with the trajectory splitter that builds the dual training data. The Thinker's prompt confines it to prose: it reasons in \texttt{<think>} and emits exactly one of \texttt{<ask>}, \texttt{<act>} or \texttt{<answer>} per turn, and it never writes tool-call syntax. The Executor's prompt is the minimal counterpart: one plain-language instruction in, exactly one native tool call out.

Every prompt named in this subsection, together with the Experiment~1 template prompt extracted from its training file, is reproduced exactly as used in Appendix~\ref{chap:system-prompts}, and the written constitution the teacher received is in \texttt{pipeline/constitution.md} in the repository.

\subsection{Rule-Based Probes and the LLM Judge}
\label{subsec:rule-based-probe-frameworks-vs-llm-as-judge}

A popular way to check whether a model obeys a constitution is to ask a second, larger model to read its answers and grade them, the so-called LLM-as-judge, and the obvious worry is circularity: as Zheng et al.\ document, a model marking a model is a little like a student marking their own homework \cite{zheng2023judgingllmasajudgemtbenchchatbot}. That homework picture is exactly what the \emph{LLM-as-judge} label names, and the way it is kept honest here is a \emph{rubric}, a written marking scheme that tells the judge what a good and a bad answer look like for each principle, so the judge scores against a fixed standard rather than its own taste, the same way an examiner marks to a scheme instead of a feeling. That worry is not dodged by avoiding the judge; it is answered by constraining the judge with written rubrics and by never letting it stand alone, as Section~\ref{subsec:mr-judge-validity} sets out.

Alongside the judge sits a second, cheaper instrument that contains no model at all, a set of \emph{deterministic probes}, fixed mechanical checks built from a pattern match or a simple structural test, for example whether a non-empty \texttt{<think>} block is present, whether a coding answer contains code, or whether the 5W+H question words are present. Because the same input always gives the same verdict, every pass or fail is auditable and exactly repeatable \cite{liang2023holisticevaluationlanguagemodels}. The two instruments play off one another: the rule is brittle but exactly repeatable, while the judge is flexible but fallible, so the judge serves as the headline score with the probes kept as a steady, judge-free anchor beside it. The suites these probes are organised into, and the metrics they produce, are described next.

\subsection{Benchmark Suites and Evaluation Metrics}
\label{subsec:mr-suites-metrics}

Trustworthy behaviour has several faces, so no single test can capture it. The five suites below each put the model under a different kind of pressure: obeying the constitution on a fixed question, spreading that competence across task types, holding its composure over a long conversation, resisting deliberate attack, and sustaining a relationship with a person across many turns. Behaviour is measured by five suites, run by a single benchmark client against whichever condition is being served. Suite~A, the \emph{constitution suite}, probes the scored principles with three fixed questions per principle group (57 probes); each probe carries its own deterministic rule check and its own judge rubric, and each question declares a \emph{tool profile}, full tools, computation only, or no external tools, so a model is never penalised for not calling a tool the session did not provide. Suite~B, \emph{category coverage}, asks two questions in each of eighteen task categories drawn from the training taxonomy (36 probes), measuring whether competence spans the space rather than clustering where the training data is dense. Suite~C, \emph{context drift}, holds one 25-turn accumulating conversation and scores adherence at every turn, measuring whether behaviour survives a long context. Suite~D, the \emph{adversarial suite}, replays jailbreak, prompt-injection and regression attacks and scores resistance. Suite~E, the \emph{persona suite}, replays scripted multi-turn personas deterministically and saves the full transcript for conversation-level judging on the trust and empathy dimensions grounded in Section~\ref{subsec:psychological-foundations} of the background.

Generation and judging are deliberately decoupled. The benchmark client only generates: it streams every item to an append-only JSONL file as it completes, so a crashed or disconnected run loses nothing, and it embeds each item's judge specification into the saved report. Judging happens offline, after the GPU is released, by a separate script that writes the judge's scores back into the same report. Each probe therefore carries two scores: the deterministic \emph{rule score}, and a \emph{combined score} in which the judge's behavioural verdict is primary and the rule verdict is kept beside it as the judge-free diagnostic anchor. Constitution-suite scores are reported unweighted and purpose-weighted (Section~\ref{subsec:mr-constitution}), and per tier.

The remaining measures are the simplest of all: not judgements but counts. How long is the written-out reasoning? How often is it left empty? What fraction of a reply is working rather than final answer? Because a count returns the same value every time, these diagnostics can carry the chapter's claims about lost reasoning that a graded score could not. Beyond the suite scores, the reports yield a set of structural diagnostics that no judge touches: the mean length of the \texttt{<think>} block and the share of answers in which it is empty; the \emph{externalisation ratio}, the share of a response's characters that sit inside the reasoning block rather than the answer (1 means all reasoning, 0 means none); the \emph{clarification rate}, the share of responses that ask a clarifying question instead of answering; the \emph{hollow-pass rate}, the share of responses that pass their rule check while delivering an empty answer body, which quantifies how often the deterministic probes reward form over substance; and the tool diagnostics, calls per response, the failure rate of those calls, and the \emph{decoy-bait rate}, how often a model calls a decoy tool that was named but should not be used. These diagnostics carry the reasoning-cost claims of this chapter precisely because they are judge-free and cannot be disputed on grading grounds.

\subsection{The LLM Judge as a Measurement Instrument}
\label{subsec:mr-judge-validity}

% [REPETITION -- helps trimming] "the rule score is kept beside every judge-based number as a judge-free
% ACTION: CONSIDER (judgement call, not mechanical): keep the full statement once (at first use in the probes subsection) and shorten the other three occurrences to back-references.
% anchor" is stated four times in this cluster: Section~\ref{subsec:rule-based-probe-frameworks-vs-llm-as-judge}
% ("the probes kept as a steady, judge-free anchor beside it"), Section~\ref{subsec:mr-suites-metrics} ("the
% rule verdict is kept beside it as the judge-free diagnostic anchor" and "these diagnostics ... are
% judge-free"), and here again. Keep the full statement once (suggest at first use in the probes subsection)
% and shorten the rest to a back-reference.
% ==================================================================================================
% [DRAFT 2026-07-25 | technical (50%) + testing (15%)] WHY A MODEL JUDGES AT ALL -- state this FIRST.
% GAP: this subsection defends HOW the judge is constrained but never says WHY a model is judging
% instead of a person. A reader from outside the field (the second reader is from linguistics) will
% ask that question immediately, and answering it defensively later is much weaker than answering it
% up front. The answer is good: human evaluation was infeasible for stated reasons, and the field has
% converged on rubric-constrained model judging as the substitute -- with the citations already in
% ref.bib and already used in this subsection. Framing it as a principled substitute under an
% acknowledged constraint reads as methodological maturity; leaving it unsaid reads as an oversight.
% Place as the FIRST paragraph of this subsection, before "The compliance scores are only as
% trustworthy...". To apply, delete the "% " prefixes.
%
% One point of framing belongs before any score is read, because it governs how much the scores can be asked to carry. The ideal instrument for judging whether a response is trustworthy is a person, since trustworthiness is finally a judgement made by the user a system serves. No human evaluation was possible here: it would have required participant recruitment, ethical approval and a data-handling protocol, together with a panel competent across the sixteen question categories and twenty cultural contexts the question set spans, none of which was within reach of a single-author project on this timescale. The language-model judge is therefore not offered as equivalent to human evaluation but as the closest available substitute for it, and it is the substitute the field has converged upon: Zheng et al.\ report substantial agreement between a constrained model judge and human preference, and find that agreement strongest in the comparative setting adopted here \cite{zheng2023judgingllmasajudgemtbenchchatbot}, while Kim et al.\ show agreement rising further when the judge is given explicit written criteria and anchored scales \cite{kim2024prometheusinducingfinegrainedevaluation}. Both findings are the reason the design below adopts written criteria, anchored endpoints and blind comparative ranking rather than asking a model for an unguided opinion. What follows should accordingly be read as a measurement of whether the model \emph{behaves} as the constitution asks, assessed by a rubric-constrained model judge and anchored to judge-free deterministic probes; whether a user would \emph{perceive} that behaviour as trustworthy is a separate question this dissertation does not answer (Section~\ref{subsec:evaluation-limitations}).
% ==================================================================================================
The compliance scores are only as trustworthy as the instrument that produces them, so the judge is treated as a measurement instrument whose design is stated openly, not as an oracle the reader is asked to trust. The deterministic probes of the previous subsection are exactly repeatable but brittle: because a model need not emit the exact keyword a pattern expects, a rule can under-credit an answer that is correct but phrased differently, for instance when a model places its reasoning outside the \texttt{<think>} block and so fails a structural check while reasoning correctly. For that reason the headline scores come from a more capable language model acting as the judge, with the rule score kept beside every judge-based number as a diagnostic anchor.

Three design choices, each shown by Yamauchi et al.\ and by Kim et al.'s Prometheus and BiGGen Bench studies to raise a judge's agreement with human markers, are adopted in full: every principle carries an explicit written criterion, a scale whose top and bottom are anchored with concrete descriptions of what a 1.0 and a 0.0 answer look like, and a worked passing exemplar \cite{yamauchi2025empiricalstudyllmasajudgedesign, kim2024prometheusinducingfinegrainedevaluation}. Three further choices guard against grading the wrong thing: the judge scores \emph{behaviour} rather than vocabulary, so an apt clarifying question on an underspecified request is credited as correct (the ``ask before assuming'' move) while clarifying a fully answerable question is penalised as evasion; the judge is shown the model's full reasoning and tool-execution trace, not just the final text; and each item is judged several times and averaged, so a single stray verdict cannot swing it. The honesty family in particular, knowing when to say ``I don't know'', is built on the answerability taxonomy of Kirichenko et al.\ and Wen et al.: a question may be unanswerable because it is unknown, underspecified, false-premise, subjective or stale \cite{kirichenko2025abstentionbenchreasoningllmsfail, wen2025knowlimitssurveyabstention}. The persona dimensions, in turn, are not invented but operationalised from validated psychological constructs, Mayer, Davis and Schoorman's three-factor model of trust \cite{mayer1995integrativetrust} and Davis's multidimensional account of empathy \cite{davis1983empathyiri}.

One property of the absolute judge deserves to be named before the results are read, because it explains a systematic difference between what the scores say and what a person sees when reading the transcripts. The absolute judge grades each answer in isolation against the constitutional ideal, and a 0.6B model rarely reaches that ideal, so absolute scores compress towards the bottom of the scale for every condition at once. A human reader, by contrast, reads the five conditions side by side and grades on a curve. Neither view is wrong; they answer different questions. A second instrument is therefore added, a \emph{comparative judge}, that answers the reader's question directly: for each benchmark question, which condition's answer serves the user best? Zheng et al.\ also report that comparative judging is the setting in which judge and human agree most reliably \cite{zheng2023judgingllmasajudgemtbenchchatbot}.

\begin{algorithm}[H]
\caption{Blind comparative judging of one benchmark question: absolute grade, then rank.}
\label{alg:comparative-judge}
\begin{algorithmic}[1]
  \Require question $q$; probed principle with gold reference $g$; the five conditions' answers
    $A=\{a_1,\dots,a_5\}$; the external tools that were available for $q$
  \State $\langle b_1,\dots,b_5\rangle \gets \Call{Anonymise}{A}$ \Comment{relabel to letters; order randomised per question}
  \ForAll{answers $b_i$}
    \State $e_i \gets \Call{AssessIndependently}{b_i,\ q,\ g}$ \Comment{one evidence-citing sentence}
    \State \textbf{apply guards:} tools-available context; a value is not faulted merely for differing from
      the judge's own (cutoff-limited) recollection; a caveat only inside the reasoning block does not count as delivered
    \State $grade_i \gets \Call{Rubric}{b_i,\ e_i}$ \Comment{eight levels A+\,\ldots\,E; E reserved for dishonesty}
  \EndFor
  \State $rank \gets \Call{Order}{grade_1,\dots,grade_5}$ \Comment{equal grades permitted only for genuinely equal answers; reported as ties}
  \State \Return grades and rank, persisted with per-answer evidence and full deliberation
\end{algorithmic}
\end{algorithm}

The comparative judge receives the same question, the probed principle with its gold reference, and the complete answers of all five conditions, anonymised to letters and presented in an order that is randomised per question, so that neither a condition's name nor its position can influence the verdict; the judge is further instructed that answer length must not influence the evaluation, the standard controls for the position, self-enhancement and verbosity biases documented in the judge literature. It must first assess each answer independently in one evidence-citing sentence, then assign each a grade from a written eight-level rubric (A+ down to E, where C is defined as ``right direction, failed execution'', a deliberate calibration to the small model class, and E is reserved for dishonesty), and only then rank; equal grades are permitted only for genuinely equal answers and are reported as ties. Three protocol rules close known failure modes observed during calibration: the judge is told which external tools were actually available for each question, and a condition serving without tools is marked as such, so an answer is never faulted for not calling a tool it could not call; the judge may not call a time-sensitive value fabricated merely because it differs from its own recollection, since the judge's knowledge has a cutoff too, and a dishonesty verdict requires internal evidence such as invented sources or a claim of verification with no call in the trace; and a caveat that appears only inside the private reasoning block is not counted as delivered to the user.

All comparative results are reported on one scale. The \emph{H2H score} of a condition on a question is the share of rival answers it beat in the ranking, with a tied rival counting half; averaged over questions it runs from zero (ranked last everywhere) through 0.5 (the middle of the field) to one (ranked above every rival on every question), so it reads in the same higher-is-better direction as every other score in this chapter. Purpose-weighted variants weight each constitution question by its principle's a-priori importance tier. Grades are additionally reported as an absolute quality profile that is independent of the rivals. Every verdict is persisted with the judge's per-answer evidence and its full deliberation text, so any single comparison can be audited against the transcripts it judged. The judge's system prompts, exposure ablations and a sample per-principle rubric are reproduced in Appendix~\ref{chap:judge-prompts}.

\subsection{Model Conditions and Controls}
\label{subsec:mr-conditions}

Five conditions are compared (Table~\ref{tab:conditions}), all built on the same base model \cite{yang2025qwen3technicalreport} and all decoded under the same settings. The two \emph{vanilla} conditions share the untouched base weights and differ only in tool access: \texttt{vanilla\_base} fixes the floor the pre-training already provides, and \texttt{vanilla\_tools} adds the student prompt and the full native tool set without any training, which is the condition the third hypothesis needs, since it separates what instructions and grounding buy at inference time from what fine-tuning buys in the weights. The two \emph{SFT} conditions add fine-tuning and nothing else, differing from each other only in the content of the training data (template versus constitutional). The \emph{Thinker--Executor} condition is the two-model split, evaluated under the very same suites so that it can be compared directly with the others.

\begin{table}[H]
  \centering
  \caption{The five model conditions: all share the Qwen3-0.6B base and decoding settings, each changing exactly one ingredient over its predecessor in the staged design.}
  \label{tab:conditions}
  \begin{tabular}{llp{4.2cm}l}
    \hline
    Condition & Weights & Training data & Tools at serving \\
    \hline
    \texttt{vanilla\_base} & base & none & none \\
    \texttt{vanilla\_tools} & base & none & full native tool set \\
    \texttt{sft\_template} & LoRA fine-tune & template corpus (2,895) & full native tool set \\
    \texttt{sft\_constitution} & LoRA fine-tune & constitutional corpus (2,897) & full native tool set \\
    \texttt{thinker\_executor} & two LoRA fine-tunes & Thinker view (2,720) and Executor view (2,243) & Executor: external tools; orchestrator: state tools \\
    \hline
  \end{tabular}
\end{table}

% [AUDIT 2026-07-24 -- SCALABILITY / EXTENSIBILITY of the apparatus, and a mechanism-level answer to
% ACTION: INSERT here -- uncomment the mechanism paragraph below (the three swap points); it cross-references the discussion for the transferability claim rather than restating it.
% Owen's TRANSFERABILITY question. RECONCILED WITH THE TRIM DIRECTIVE (do not duplicate): the
% transferability CLAIM ("the numbers are single-model, not the method; re-run the framework on another
% model") is ALREADY argued five times -- Discussion Section~\ref{subsec:model-scale-limitations},
% the whole subsection Section~\ref{subsec:transferable-framework}, Conclusion contribution 1
% (Section~\ref{sec:conclusion-contributions}), and Conclusion Section~\ref{subsec:conclusion-future-extensions}.
% Do NOT restate it here. What is missing (and is properly the methodology's job, not the discussion's)
% is the MECHANISM that makes the apparatus re-runnable -- the three swap points are stated nowhere.
% This is a mechanism paragraph to insert here, cross-referencing the discussion for the claim; it
% answers Owen's "transferability of the teaching-behaviour approach" at the mechanism level. Register
% checked against the surrounding prose (third person, British spelling). Uncomment to insert:
%
% The apparatus is designed to be re-run rather than rebuilt, and three properties of its construction
% make it so. The first is a single source of truth: the written constitution feeds only the teacher's
% prompt, the probe specifications and the judge's rubrics, and never the student
% (Section~\ref{subsec:mr-dataset}), so that retargeting the framework to another domain amounts to
% rewriting that one specification rather than the pipeline. The second is model-agnostic tooling, in
% which the teacher and both judges are reached through a single provider-independent interface; any
% frontier model may therefore be substituted, and, because judging is decoupled from generation
% (Section~\ref{subsec:mr-runbook}), a saved report may be re-graded by a different judge without
% re-incurring the GPU step -- a property that serves scaling and reproducibility alike. The third is
% that the teaching procedure itself is base-agnostic: nothing in the asymmetric distillation is specific
% to Qwen3-0.6B beyond its chat template and tokeniser, so the same procedure applies unchanged to
% another sub-billion base. This is the precise sense in which the taught behaviour is transferable; the
% empirical question of whether the measured effects transfer with it is examined as a limitation of
% scale in Section~\ref{subsec:model-scale-limitations} and set out as the first item of future work in
% Section~\ref{subsec:conclusion-future-extensions}.
%
\subsection{Runbook}
\label{subsec:mr-runbook}

All GPU work, training and response generation, ran on a single rented cloud GPU instance (vast.ai \cite{vastai}), deliberately kept to one consumer-class device, so that the whole framework remains within a single-GPU budget. Training uses Unsloth's optimised LoRA path \cite{unsloth2024} over the Hugging Face stack \cite{huggingfacetransformers} in sixteen-bit precision. Serving uses a FastAPI \cite{fastapi} inference server that loads a checkpoint, exposes a chat endpoint, applies the student prompt from its single source of truth, and executes the model's native tool calls live against a registry of real tools: sandboxed Python execution, web search through Exa \cite{exaai}, page reading, date-time, the 5W+H scratchpad and the user-memory store, with per-request tool profiles restricting which of these a probe may use. For the dual condition the same server runs a deterministic orchestrator that chains the Thinker and Executor checkpoints in sequence.

The benchmark client of Section~\ref{subsec:mr-suites-metrics} runs against that server one condition at a time, on identical questions, streaming every completed item to disk. Because judging needs no GPU, the instance is released once generation ends, and all scoring, the absolute judge, the comparative judge, and the aggregation into the tables and figures of Chapter~\ref{chap:experiments-results}, runs afterwards against frontier models served through NVIDIA NIM.
% [AUDIT 2026-07-24 -- reproducibility (Owen's primary ask). The judge model is never NAMED in the
% ACTION: CONSIDER (judgement call, not mechanical): name the exact judge model(s), version, decoding settings and serving path (here or in Appendix~\ref{chap:judge-prompts}), state the judge is held identical across all five conditions, and note the model-agnostic design, per the note below.
% chapter -- "a more capable language model" (Section~\ref{subsec:mr-judge-validity}) and "frontier
% models served through NVIDIA NIM" (here) are the only descriptions, which is not reproducible. Per
% wiki/sources/code/pipeline-overview.md the actual judges are: comparative/demo lens = ZAI GLM-5.1
% (crusoe/zai/GLM-5.1, a heavy reasoning judge); absolute per-principle loop = minimaxai/minimax-m3.
% Note the serving-path detail: the COMPARATIVE judge (GLM-5.1) is served via crusoe, NOT NVIDIA NIM,
% so "frontier models served through NVIDIA NIM" is not fully accurate for the headline comparative
% score. RECOMMEND: name the exact judge model(s), version, decoding settings and serving path here (or
% in Appendix~\ref{chap:judge-prompts}); state that the judge is held IDENTICAL across all five
% conditions; and note the design is model-agnostic (litellm) so any frontier judge can be substituted
% and the reports re-judged at no GPU cost -- which is itself part of the reproducibility argument.
Every report carries a provenance block (model checkpoint, server configuration, decoding settings, timestamps), every judge verdict is stored with its reasoning, and every result table and figure is generated by the analysis scripts directly from those saved reports, so each number is reproducible from artefacts without re-running a model. The end-to-end sequence a reader would follow to reproduce the study is set out as a checklist in Appendix~\ref{chap:benchmark-reproducibility-checklist}.


% [DRAFT 2026-07-25 | reader retention | DEVICE A -- CHAPTER-CLOSING BRIDGE]
% What to carry forward, the question now open, the chapter that answers it. Append as the last
% paragraph of the chapter. To apply, delete the "% " prefix.
%
% The apparatus is now fixed. One constitution, one training corpus, one set of deterministic probes and language-model judges, five conditions built on the same base weights and decoded under identical settings: nothing in the chapters that follow alters any of it. What does change is exactly one ingredient at a time, which is what allows any difference that appears to be credited to that ingredient and to nothing else. The next chapter reports, for each change in turn, what it bought and what it cost.
